% ============================================================
% Section 158: 双$\delta$势垒中的共振态能级
% ============================================================
\section{量子力学：双$\delta$势垒中的共振态能级}
\label{sec:double-delta-resonance}

\begin{modelbox}[题目：双$\delta$势垒中的共振态能级]
粒子在双$\delta$势垒
\[
V(x)=V_0[\delta(x+a)+\delta(x-a)]
\]
作用下运动，求共振态能级。
\end{modelbox}

\begin{tipbox}[考点快速分析]
\begin{itemize}
    \item \textbf{共振态}：内部区域波函数振幅远大于外部，粒子被``准束缚''
    \item \textbf{宇称分类}：偶宇称和奇宇称分别求解
    \item \textbf{共振条件}：外部振幅 $C\ll1$，$\cos ka\approx0$（偶）或 $\sin ka\approx0$（奇）
    \item \textbf{极限}：强势垒 $b\to0$ 时退化为宽度 $2a$ 的无限深势阱
\end{itemize}
\end{tipbox}

\begin{derivationbox}[标准解题过程]
\textbf{1. 基本设定}

令 $k=\sqrt{2mE}/\hbar$，$b=\hbar^2/(mV_0)$。在 $x\neq\pm a$ 区域方程为 $\psi''+k^2\psi=0$。

\textbf{2. 偶宇称共振态}

$\psi(-x)=\psi(x)$。内部（$|x|<a$）：$\psi(x)=\cos kx$；外部（$x>a$）：$\psi(x)=C\sin(kx+\theta)$。

$x=a$ 处边界条件：
\begin{align*}
\cos ka&=C\sin(ka+\theta),\\
Ck\cos(ka+\theta)+k\sin ka&=\frac{2}{b}\cos ka.
\end{align*}

共振条件 $C\ll1$ 要求 $|\cos ka|\ll1$，即 $ka\approx\left(n+\dfrac{1}{2}\right)\pi$。

设 $ka=\left(n+\dfrac{1}{2}\right)\pi-\varepsilon$，$|\varepsilon|\ll1$。由第二式得 $1\approx\dfrac{2ka}{b}\varepsilon$，故 $\varepsilon\approx\dfrac{kb}{2}\approx\dfrac{(2n+1)\pi b}{4a}$。

修正后的 $k$：
\[
k\approx\frac{(2n+1)\pi}{2a}\left(1-\frac{b}{2a}\right).
\]

能级：
\begin{equation}
\boxed{E_n^{\text{偶}}\approx\frac{(2n+1)^2\pi^2\hbar^2}{8ma^2}\left(1-\frac{b}{a}\right),\quad n=0,1,2,\ldots}
\end{equation}

\textbf{3. 奇宇称共振态}

$\psi(-x)=-\psi(x)$。内部：$\psi(x)=\sin kx$；外部：$\psi(x)=D\sin(kx+\varphi)$。

共振条件 $D\ll1$ 要求 $\sin ka\approx0$，即 $ka\approx n\pi$（$n=1,2,\ldots$）。

类似推导得：
\begin{equation}
\boxed{E_n^{\text{奇}}\approx\frac{n^2\pi^2\hbar^2}{2ma^2}\left(1-\frac{b}{a}\right),\quad n=1,2,\ldots}
\end{equation}

\textbf{4. 统一公式与极限}

统一写为 $E_n\approx\dfrac{n^2\pi^2\hbar^2}{8ma^2}\left(1-\dfrac{b}{a}\right)$（$n=1,2,\ldots$），其中 $n$ 为奇数对应偶宇称，$n$ 为偶数对应奇宇称。

强势垒极限（$b\to0$）：退化为宽度 $2a$ 的无限深势阱能级。

共振态存在条件：$\dfrac{b}{a}\ll\dfrac{4}{n\pi}$，即 $V_0\gg\dfrac{n\hbar^2}{2ma}$。
\end{derivationbox}

\begin{formulabox}[答案汇总]
\begin{center}
\begin{tabular}{ll}
\toprule
\textbf{结论} & \textbf{结果} \\
\midrule
偶宇称共振能级 & $E_n^{\text{偶}}\approx\dfrac{(2n+1)^2\pi^2\hbar^2}{8ma^2}\left(1-\dfrac{b}{a}\right)$ \\
奇宇称共振能级 & $E_n^{\text{奇}}\approx\dfrac{n^2\pi^2\hbar^2}{2ma^2}\left(1-\dfrac{b}{a}\right)$ \\
强势垒极限 & 宽度 $2a$ 无限深势阱能级 \\
\bottomrule
\end{tabular}
\end{center}
\end{formulabox}

\begin{insightbox}[物理图像]
\begin{itemize}
    \item 双$\delta$势垒系统没有严格的束缚态，但存在共振态（准束缚态），能量略低于同宽度的无限深势阱能级。
    \item 共振态的物理图像是：粒子在双势垒之间来回反射，形成类似 Fabry-Pérot 干涉仪的驻波。当势垒足够强时，透射概率极低，粒子被``准束缚''在内部区域。
    \item 该模型是理解量子阱、共振隧穿二极管等器件中准束缚态和共振隧穿现象的基础。
\end{itemize}
\end{insightbox}
